\chapter{Introducción}
\label{chap:introduccion}

\lettrine{L}{os} entornos virtuales en \textit{3D} han resultado ser buenos escenarios para
la investigación de algoritmos de \gls{ia}.
A diferencia de entornos de juegos clásicos como el ajedrez, los entornos virtuales en \textit{3D} ofrecen
percepción visual en diferentes perspectivas (primera persona, tercera persona, etc.),
espacios de acción continuos o de alta dimensionalidad y características temporales complejas,
aumentando mucho más su dificultad y acercándolos más al funcionamiento del mundo
real~\cite{DBLP:journals/nature/SilverHMGSDSAPL16,DBLP:journals/nature/VinyalsBCMDCCPE19}.
Plataformas como Atari~\cite{DBLP:journals/nature/MnihKSRVBGRFOPB15} o VizDoom~\cite{kempka2016vizdoom}
han sido referentes clave para el uso de aprendizaje por refuerzo y planificación autónoma
en los videojuegos.

Dentro de esta familia de entornos virtuales, \textit{Minecraft} ocupa un lugar singular.
Se trata de un videojuego \textit{3D} de mundo abierto en primera persona
centrado en la recolección de recursos, la construcción de estructuras,
la fabricación de objetos y la supervivencia~\cite{guss2019minerl}. Se caracteriza por sus mundos, 
generados de forma procedimental y formados por bloques que el jugador
puede modificar, lo que genera un espacio de acciones grande y variado.
A diferencia de otros entornos de referencia, \textit{Minecraft} carece de un objetivo 
único predefinido, lo que obliga a los agentes a planificar
horizontes temporales largos con recompensas dispersas y, a menudo,
a trabajar en situaciones nunca vistas durante el entrenamiento.

El desafío más complejo que presenta el juego consiste en derrotar al \textit{Ender Dragon},
un jefe final que requiere la planificación y ejecución de una gran red de tareas anidadas entre sí.
Tarea que aún no ha sido completada usando únicamente agentes autónomos, sin ninguna intervención humana.

A nivel de progreso, todos los objetivos que plantee el jugador pueden alcanzarse de varias formas y empleando diferentes herramientas.
A partir de esta premisa, los jugadores pueden fabricar materiales y objetos específicos que requieren la recolección previa 
de una variedad de recursos, lo que da lugar a un complejo sistema de creación organizado como un árbol jerárquico
(véase el Apéndice~\ref{apx:iron-pickaxe}).

Como ejemplo, si el jugador quisiera un pico de hierro debería: recolectar primero madera rompiendo un árbol, luego fabricar una mesa de trabajo para 
fabricar nuevos objetos; dentro de ella, fabricar un pico de madera para recolectar piedra y, a continuación, un pico de piedra 
con el que extraer carbón y mineral de hierro; después, construir un horno para fundir el mineral y obtener lingotes de hierro para, 
finalmente, fabricar el pico de hierro. Además, el jugador debe encontrar en el mundo todos los recursos necesarios.

Esta tarea, que puede parecer sencilla para un jugador humano, supone un gran desafío para
los algoritmos de \gls{ia}, ya que tanto el entorno como el espacio de acciones son muy grandes y complejos,
y presentan además un fuerte componente temporal difícil de capturar.

Ante este desafío se pueden aplicar paradigmas de \gls{ia} muy diferentes entre sí,
desde la planificación simbólica basada en conocimiento experto
hasta el aprendizaje a partir de imitación o usando la propia interacción con el entorno.
Este trabajo, por tanto, compara estas tres familias: planificación jerárquica de tareas (\gls{htn}),
aprendizaje por imitación (\gls{il}) y aprendizaje por refuerzo (\gls{rl}), sobre una misma tarea
en \textit{Minecraft} y bajo idénticas condiciones de evaluación, con el fin de analizar
las fortalezas y limitaciones de cada enfoque y 
la relación entre conocimiento experto, datos y coste de cómputo.

\section{Objetivos}
El trabajo plantea las siguientes metas finales:
\begin{itemize}
    \item Diseñar e implementar agentes autónomos basados en distintos paradigmas de \gls{ia}. Específicamente \gls{htn}, \gls{il} y \gls{rl}, enfocados en la resolución de tareas en \textit{Minecraft}.
    \item Desarrollar, en caso necesario, conjuntos de datos específicos que permitan el entrenamiento de los modelos seleccionados.
    \item Definir métricas adecuadas para la evaluación del rendimiento de los agentes en el entorno.
    \item Analizar y comparar el rendimiento, eficacia y limitaciones de cada paradigma.
    \item Combinar sistemas de forma eficiente, utilizando diferentes lenguajes de programación para el desarrollo de modelos 
    y el control de agentes, así como la comunicación entre distintos entornos de desarrollo.
\end{itemize}

\section{Estructura de la Memoria}
La organización de este trabajo es la siguiente: tras el
Capítulo~\ref{chap:introduccion} de introducción, se hace una revisión del estado del arte en el
Capítulo~\ref{chap:sota}, enfocado en técnicas de \gls{ia} aplicadas a videojuegos y, 
en particular, \textit{Minecraft}. A continuación, el
Capítulo~\ref{chap:marco_teorico} introduce los fundamentos teóricos de las tres familias
de técnicas comparadas (\gls{htn}, \gls{il} y \gls{rl}), mientras que el Capítulo~\ref{chap:metodologia} detalla la metodología de
desarrollo, las fases del proyecto, su planificación y los conjuntos de datos empleados. El
Capítulo~\ref{chap:arquitectura} describe la arquitectura del sistema común a todas las
iteraciones. Sobre esta arquitectura se construyen las tres iteraciones que forman el núcleo del
trabajo: el Capítulo~\ref{chap:iter1} introduce al agente \gls{htn}, que actúa de experto;
el Capítulo~\ref{chap:iter2}, el agente \gls{il} que aprende clonando a dicho experto;
y el Capítulo~\ref{chap:iter3}, el agente \gls{rl} que aprende desde cero por interacción
con el entorno. Tras esto, en el Capítulo~\ref{chap:discusion-resultados} se comparan las tres
técnicas bajo una métrica común y se analiza la relación entre conocimiento experto, datos
y coste de cómputo. Por último, se dedica el Capítulo~\ref{chap:conclusions}
a exponer las conclusiones y las posibles líneas futuras de investigación
obtenidas del análisis realizado en el trabajo.